\section{PIOPs over fields for cyclotomic ring arithmetic}
\label{sec:ring-piop}

We present a (multilinear) PIOP for proving matrix-vector products $As \equiv t$ in power-of-two cyclotomic rings $\ring_p = \ZZ_p[X]/(X^{2^\nu}+1)$ with prime moduli $p$, where all messages and polynomial coefficients are in $\extensionfield$ for a sufficiently large extension field $\extensionfield/\field_p$. This is given in~\cref{sec:matrix-vector-piop}. Subsequent sections present multiple variants and optimizations of this protocol, as we now describe.

\paragraph{Scope and restrictions.}
Our protocol requires three structural assumptions on the ring $\ring_p$: the modulus $p$ is prime, the cyclotomic polynomial is $\Phi_{2n}(X) = X^n + 1$ with $n = 2^\nu$ a power of two, and $p$ splits $\Phi_{2n}$ completely (i.e., $p \equiv 1 \pmod{2n}$, so that $X^n + 1$ factors into $n$ distinct linear factors over $\field_p$). Each assumption serves a distinct purpose:
\begin{itemize}
    \item \emph{Prime modulus.} The sumcheck protocol, the polynomial commitment scheme, and the extension-field arithmetic all require that we work over a field. A composite modulus $q$ is not a field, so one cannot run the sumcheck or PCS directly over $\ZZ_q$. The standard workaround is to factor $q = \prod p_i$ and verify the relation modulo each prime factor independently; our protocol handles each such factor.
    \item \emph{Fully splitting.} The CRT-domain approach requires that $X^n + 1$ admits $n$ distinct roots $\omega, \omega^3, \ldots, \omega^{2n-1}$ in $\field_p$, where $\omega$ is a primitive $2n$-th root of unity. This full splitting is what makes the ring product become a pointwise product after the CRT transform, yielding the key identity $t_C(i,k) = \sum_j A_C(i,j,k)\cdot s_C(j,k)$ that Sumcheck~1 verifies. When $\Phi_{2n}$ does not split completely over $\field_p$ (i.e., $p \not\equiv 1 \pmod{2n}$), the roots of $X^n + 1$ live in a proper extension of $\field_p$, and the CRT gives evaluation over that extension rather than over $\field_p$ itself; the protocol structure could in principle be adapted to this setting at the cost of increasing the base field size and the CRT matrix complexity, but we do not pursue this here.
    \item \emph{Power-of-two degree.} This is primarily an efficiency assumption rather than a correctness requirement. The power-of-two structure $n = 2^\nu$ gives the CRT matrix $C_{k,\ell} = \omega^{(2k+1)\ell}$ a recursive dyadic periodicity: each factor $f_j(k) = (1-\gamma_j) + \gamma_j \omega^{(2k+1)2^j}$ has period $2^{\nu-j}$ in $k$, which is what enables the $O(n)$-time verifier evaluation of~$\tilde{C}$ in \cref{lem:m-eval}. For a general cyclotomic $\Phi_m(X)$ of degree $\varphi(m)$, the roots do not exhibit this clean doubling structure, so the factored product-table construction would not apply and the verifier's $\tilde{C}$ evaluation would cost $O(\varphi(m)^2)$ in the worst case. The power-of-two case also aligns with the DP24 packing reduction (\cref{sec:dp24-packing}), which uses a tower of quadratic extensions $\field_p \subset \field_{p^2} \subset \cdots \subset \extensionfield$ whose depth is $\log_2[\extensionfield:\field_p]$.
\end{itemize}
These restrictions are broadly compatible with deployed lattice-based cryptography: NIST standards (ML-KEM, ML-DSA) and most lattice-based signature schemes use power-of-two cyclotomic rings with primes chosen so that $p \equiv 1 \pmod{2n}$, and composite moduli decompose into such primes.

\paragraph{Why matrix-vector products.}
The relation $t = As$ over $\ring_q$ is the common algebraic core of lattice-based cryptography: MLWE/RLWE encryption, Dilithium/Falcon/Hawk signatures, FHE operation proofs, Ajtai hash commitments, and lattice-based VRFs all reduce --- individually or in batches --- to one or more instances of this relation. Matrix-matrix products $T = AS$ reduce to matrix-vector products via column batching: the verifier samples $\rho \randpick \extensionfield^{M'}$ and the single relation $T\rho = A(S\rho)$ is checked, while multiple independent instances $A s_i = t_i$ batch to a single inner product as described in \cref{par:batch-verification}. We therefore present the protocol in full generality for $A \in \ring_q^{N \times M}$, then give optimized treatments for the inner product case $N = 1$ (which covers RSIS/SWIFFT, Falcon verification, and all batched multi-instance settings).

Beyond lattice cryptography, the matrix-vector relation over cyclotomic rings is a natural building block for proving more general NP statements that involve ring arithmetic. Any computation expressed as an intermediate representation over $\ring_q$ --- such as an arithmetic circuit whose gates are ring additions and multiplications, or a sequence of NTT-based operations --- decomposes into a series of bilinear relations over $\ring_q$, each of which is an instance of matrix-vector multiplication (with $N = M = 1$ for scalar multiplication, $N = 1$ for inner products, and general $N, M$ for linear maps). Because our protocol is a sumcheck-based PIOP, it composes naturally with other sumcheck protocols: claims produced by one sumcheck invocation can feed directly into another, enabling modular stacking without additional commitment overhead. The CRT-domain structure we exploit is also not specific to matrix-vector products; any relation that becomes pointwise after applying the CRT transform admits the same two-sumcheck architecture (a pointwise-check sumcheck followed by a CRT consistency sumcheck). While we do not pursue a fully general ring-arithmetic compiler in this work, the techniques developed here --- CRT-domain verification, holographic indexing, packing, and oracle reduction --- provide the essential ingredients for doing so.

For $n = 2^\nu$ and a $N \times M$ matrix $A$, our protocol has prover complexity $O(NMn)$\psiv{check this}, verifier complexity $O(n+\log NM)$, and proof size $O(\log Mn)$. The additive $O(n)$ term in verifier complexity arises from a single evaluation of the CRT matrix's multilinear extension (\cref{lem:m-eval}) and cannot be eliminated without fundamentally changing the protocol structure. An alternative ``coefficient-domain'' approach based on ring switching and a convolution sumcheck can in principle replace this $O(n)$ with $O(\log n)$, but it produces cross-polynomial product constraints that are incompatible with multiple of the optimizations our protocol benefits from. In practice, $n$ is often small (e.g., $n = 64$ for SWIFFT) and the constant factor in the $O(n)$ evaluation is small, so it does not dominate concrete verifier performance. \psiv{Give example estimate of how much it would contribute to total verifier computation.}
%That said, in~\cref{sec:logger-verifier} we introduce a $O(\log^2 n+\log NM)$ verifier variant, with the tradeoff of it having $O(\log M + \log^2 n)$ proof size.

We write the main PIOP construction in~\cref{sec:matrix-vector-piop} in the ``commit and prove'' style, where the verifier begins with three oracles to polynomial representations of matrix $A$ and vectors $s$ and $t$ (that we assume were earlier sent by the prover). When $A$ is constant (e.g., proving knowledge of an MSIS opening) and/or $t$ is public, we find ourselves in the holographic setting where an indexer can output variations of these oracles that allow for a more efficient argument. This is explained in~\cref{sec:holographic-piop}.

\subsection{A PIOP for cyclotomic matrix-vector multiplication}
\label{sec:matrix-vector-piop}

Let $\omega \in \field_p$ be a primitive $2n$-th root of unity and let $\tilde{C} \in \field_p^{\multilin}[\bm{K},\bm{L}]$ denote the multilinear extension of the CRT matrix $C = \crtnosubscript_{n,p}$,
\[
    C_{k,\ell} = \omega^{(2k+1)\ell}, \qquad k,\ell \in \{0,1\}^\nu,
\]
which maps coefficient vectors to evaluations at the roots of $X^n + 1$. Define virtual polynomials
\[
    \begin{array}{ll}
    A_C(i,j,k) &:= \sum_{\ell \in \bits^\nu} \tilde{C}(k,\ell)\cdot\tilde{A}(i,j,\ell), \\
    s_C(j,k) &:= \sum_{\ell \in \bits^\nu} \tilde{C}(k,\ell)\cdot\tilde{s}(j,\ell), \\
    t_C(i,k) &:= \sum_{\ell \in \bits^\nu} \tilde{C}(k,\ell)\cdot\tilde{t}(i,\ell).
    \end{array}
\]
When $t \equiv A s \mod p$, the CRT transform gives the pointwise relation $t_C(i,k) = \sum_{j \in \bits^\mu} A_{C}(i,j,k)\cdot s_{C}(j,k)$ for all $(i,k) \in \bits^{\tau+\nu}$.

\begin{figure}[H]
\caption{Protocol $\Pi_{\mathsf{RingMul}}=(\prover, \verifier)$: a PIOP for checking matrix-vector multiplication in a cyclotomic ring $\ring_p = \ZZ_p[X]/(X^{2^\nu}+1)$ over $\field_p$}
\label{prot:ringmul}
\begin{algorithm}[H]
\textbf{Input:} The verifier receives oracles $\oracle{\tilde{A}} \in \field_{p}^{\multilin}[\bm{I},\bm{J},\bm{K}]$, $\oracle{\tilde{s}} \in \field_{p}^{\multilin}[\bm{J},\bm{K}],$ and $\oracle{\tilde{t}} \in \field_{p}^{\multilin}[\bm{I},\bm{K}]$. The prover receives matrix $A \in \ring_{p}^{N\times M}$ and vectors $s \in \ring_{p}^{M}$, $t \in \ring_{p}^{N}$ such that
\[
    \tilde{A}(i,j,k) = \coeff_k(A_{i,j}), \qquad
    \tilde{s}(j,k) = \coeff_k(s_j), \qquad
    \tilde{t}(i,k) = \coeff_k(t_i)
\]
holds for $i \in \bits^\tau, j \in \bits^\mu, k \in \bits^\nu$, where $\tau = \log N$, $\mu = \log M$, $\nu = \log n$, and $n = 2^\nu$ is the degree of the cyclotomic polynomial.

\begin{algorithmic}[1]
    \STATE The verifier samples $\alpha \randpick \extensionfield^{\tau}$, $\alpha' \randpick \extensionfield^{\nu}$.

    \STATE The prover sends $e_t$ claimed to equal $t_C(\alpha, \alpha')$.

    \STATE $\prover$ and $\verifier$ engage in a sumcheck for the claim
    \[
        e_t = \sum_{j \in \bits^{\mu}} A_C(\alpha, j, \alpha') \cdot s_C(j, \alpha')
    \]
    with challenges $\beta \randpick \extensionfield^{\mu}$, reducing to the claim $\sigma' = A_C(\alpha, \beta, \alpha') \cdot s_C(\beta, \alpha')$.

    \STATE The prover sends $e_A$ claimed to equal $A_C(\alpha, \beta, \alpha')$. The verifier computes $e_s \gets \sigma' / e_A$ and checks $e_A \neq 0$.\footnote{The event $A_C(\alpha, \beta, \alpha') = 0$ occurs with negligible probability at most $(\tau + \mu + \nu)/|\extensionfield|$.}

    \STATE The verifier samples batching challenge $\xi \randpick \extensionfield$.

    \STATE $\prover$ and $\verifier$ engage in a sumcheck for the claim
    \[
        e_t + \xi \cdot e_A + \xi^{2} \cdot e_s  = \sum_{\ell \in \bits^{\nu}} \tilde{C}(\alpha', \ell) \cdot \Big[\tilde{t}(\alpha, \ell) + \xi \cdot \tilde{A}(\alpha, \beta, \ell) + \xi^{2} \cdot \tilde{s}(\beta, \ell)\Big]
    \]
    with folding challenges $\gamma \randpick \extensionfield^{\nu}$, reducing to the claim
    \[
        \sigma'' = \tilde{C}(\alpha', \gamma) \cdot \Big[\tilde{t}(\alpha, \gamma) + \xi \cdot \tilde{A}(\alpha, \beta, \gamma) + \xi^{2} \cdot \tilde{s}(\beta, \gamma)\Big].
    \]

    \STATE The verifier evaluates $\tilde{C}(\alpha', \gamma)$ in $O(n)$ field operations (\cref{lem:m-eval}), queries $e'_t \gets \oracle{\tilde{t}(\alpha, \gamma)}$, $e'_A \gets \oracle{\tilde{A}(\alpha, \beta, \gamma)}$, $e'_s \gets \oracle{\tilde{s}(\beta, \gamma)}$, and checks
    \[
        \sigma'' = \tilde{C}(\alpha', \gamma) \cdot \big(e'_t + \xi \cdot e'_A + \xi^{2} \cdot e'_s\big).
    \]
\end{algorithmic}
\end{algorithm}
\end{figure}

%%%%%%%%%%%%%%%%%%%%%%%%%%%%%%%%%%%%%%%%%%%%%%%%%%%%%%%%%%%%%%%%%%%%%%%%%%%%%%%
%%%%%%%%%%%%%%%%%%%%%%%%%%%%%%%%%%%%%%%%%%%%%%%%%%%%%%%%%%%%%%%%%%%%%%%%%%%%%%%
\subsection{Completness, round-by-round knowledge soundness, and efficiency}

\begin{lemma}\label{lem:m-eval}
    Let $\omega \in \field_p$ be a primitive $2n$-th root of unity with $n = 2^\nu$, and let $\tilde{C} \in \field_p^{\mathsf{ml}}[\bm{K},\bm{L}]$ denote the multilinear extension of $C = \crtnosubscript_{n,p}$, given by $C_{k,\ell} = \omega^{(2k+1)\ell}$ for $k,\ell \in \bits^\nu$. Then for any $\alpha', \gamma \in \extensionfield^\nu$, the evaluation $\tilde{C}(\alpha', \gamma)$ can be computed in $O(n)$ field operations.
\end{lemma}

\begin{proof}
    We first derive a factored representation by separating the sum over~$\ell$.
    Since $\omega^{(2k+1)\ell} = \prod_{j=0}^{\nu-1} \omega^{(2k+1)\ell_j 2^j}$
    and $\meq(\gamma,\ell) = \prod_{j} \bigl((1-\gamma_j)(1-\ell_j) + \gamma_j \ell_j\bigr)$,
    the inner sum over $\ell$ factors across coordinates:
    \begin{align}
        \sum_{\ell \in \bits^\nu} \omega^{(2k+1)\ell} \, \meq(\gamma,\ell)
        &\;=\; \prod_{j=0}^{\nu-1}\; \sum_{\ell_j \in \{0,1\}} \omega^{(2k+1)\ell_j 2^j}\, \bigl((1-\gamma_j)(1-\ell_j) + \gamma_j \ell_j\bigr) \notag \\
        &\;=\; \prod_{j=0}^{\nu-1} \Big[(1-\gamma_j) + \gamma_j\, \omega^{(2k+1)2^j}\Big]. \label{eq:Mfactored}
    \end{align}
    Define the univariate polynomial
    \[
        g(X) \;:=\; \prod_{j=0}^{\nu-1} \Big[(1-\gamma_j) + \gamma_j\, X^{2^j}\Big],
    \]
    which has degree $\sum_{j=0}^{\nu-1} 2^j = n - 1$. By~\eqref{eq:Mfactored}, $g(\omega^{2k+1}) = \prod_j [(1-\gamma_j) + \gamma_j\, \omega^{(2k+1)2^j}]$, so
    \begin{equation}\label{eq:Mfactored2}
        \tilde{C}(\alpha', \gamma) \;=\; \sum_{k \in \bits^\nu} \meq(\alpha', k)\cdot g(\omega^{2k+1}).
    \end{equation}

    To evaluate the right-hand side in $O(n)$ time, we exploit the periodicity of each factor of~$g$.
    Write $g(\omega^{2k+1}) = \prod_{j=0}^{\nu-1} f_j(k)$ where
    \[
        f_j(k) \;:=\; (1-\gamma_j) + \gamma_j\, \omega^{(2k+1)2^j}.
    \]
    Since $\omega$ has order $2n = 2^{\nu+1}$, the element $\omega^{2^{j+1}}$ has order $2^{\nu-j}$, and therefore $\omega^{(2k+1)2^j} = \omega^{2^j} \cdot (\omega^{2^{j+1}})^k$ is periodic in~$k$ with period $2^{\nu-j}$.
    In particular:
    \begin{itemize}
        \item $f_{\nu-1}$ has period $2$ in $k$ \,($\omega^{2^\nu} = -1$, so $f_{\nu-1}(k)$ depends only on $k \bmod 2$),
        \item $f_{\nu-2}$ has period $4$, \;\ldots, \; $f_0$ has period $n$.
    \end{itemize}

    We compute the product table $[g(\omega^{2k+1})]_{k=0}^{n-1}$ iteratively.
    Initialize with the partial product $P_{\nu-1}(k) := f_{\nu-1}(k)$ for $k \in \{0, 1\}$, at a cost of $O(1)$.
    For $s = \nu-2, \nu-3, \ldots, 0$: extend $P_{s+1}$ from its period $2^{\nu-s-1}$ to period $2^{\nu-s}$ by replication, and compute
    $P_s(k) := P_{s+1}(k) \cdot f_s(k)$ for $k = 0,\ldots, 2^{\nu-s}-1$.
    The cost at stage~$s$ is $O(2^{\nu-s})$ multiplications, so the total is
    \[
        \sum_{s=0}^{\nu-1} 2^{\nu-s} \;=\; 2 + 4 + \cdots + 2^\nu \;=\; 2(2^\nu - 1) \;=\; O(n).
    \]
    After the final stage, $P_0(k) = g(\omega^{2k+1})$ for all $k \in \{0,\ldots,n-1\}$.
    The table $[\meq(\alpha',k)]_{k \in \bits^\nu}$ is computed in $O(n)$ by the standard recursive formula.
    The evaluation~\eqref{eq:Mfactored2} is then an inner product of two length-$n$ vectors, costing $O(n)$.
\end{proof}

\subsection{A holographic variant}
\label{sec:holographic-piop}

When the matrix $A$ is public (e.g., given by the structure of a hash function such as SWIFFT, or by a lattice scheme's public key), the oracle $\oracle{\tilde{A}}$ need not be sent by the prover: it can be provided by an indexer in a preprocessing phase. We now describe how to exploit this by committing directly to the CRT-domain representation $A_C$, which eliminates the CRT consistency check for~$A$ from Sumcheck~2 entirely.

\paragraph{Indexed CRT oracle.}
Recall that $A_C(i,j,k) = \sum_{\ell \in \bits^\nu} \tilde{C}(k,\ell)\cdot\tilde{A}(i,j,\ell)$. Since $A$ is public, the indexer can compute $A_C$ explicitly and provide the verifier with an oracle $\oracle{A_C} \in \extensionfield^{\multilin}[\bm{I},\bm{J},\bm{K}]$ in a preprocessing phase. The crucial observation is that the verifier can now query $A_C(\alpha, \beta, \alpha')$ directly at Step~4 of \cref{prot:ringmul}, obtaining $e_A$ without relying on the prover. This removes $\tilde{A}$ from the batched sum in Sumcheck~2, which becomes
\[
    e_t + \xi \cdot e_s = \sum_{\ell \in \bits^{\nu}} \tilde{C}(\alpha', \ell) \cdot \Big[\tilde{t}(\alpha, \ell) + \xi \cdot \tilde{s}(\beta, \ell)\Big],
\]
with folding challenges $\gamma \randpick \extensionfield^\nu$ reducing to the final claim
\[
    \sigma'' = \tilde{C}(\alpha', \gamma)\cdot\big(e'_t + \xi \cdot e'_s\big),
\]
verified by querying $e'_t \gets \oracle{\tilde{t}(\alpha,\gamma)}$ and $e'_s \gets \oracle{\tilde{s}(\beta,\gamma)}$. The oracle $\oracle{\tilde{A}}$ no longer appears.

\paragraph{Additional simplification when $t$ is public.}
In many lattice-based settings, $t$ is also public --- for instance, in Module-SIS where $t = 0$, or in Module-LWE where $t$ is the public ciphertext. In this case, an indexer can similarly provide $\oracle{t_C}$, and the verifier obtains $e_t$ by querying $t_C(\alpha, \alpha')$ directly. Sumcheck~2 then reduces to a CRT consistency check on $\tilde{s}$ alone:
\[
    e_s = \sum_{\ell \in \bits^{\nu}} \tilde{C}(\alpha', \ell) \cdot \tilde{s}(\beta, \ell).
\]
This is the minimal form of the protocol: a single prover-committed oracle $\oracle{\tilde{s}}$, two sumchecks, and one PCS evaluation.

\paragraph{Prover and indexer preprocessing.}
The indexer computes $A_C$ by applying the CRT transform (an $n$-point NTT at the odd powers of $\omega$) to each of the $NM$ entries of~$A$, for a total cost of $O(NMn\log n)$ field operations, or $O(NMn)$ using a batched radix-$2$ NTT. If $t$ is also indexed, the same applies at cost $O(Nn \log n)$. On the prover side, when $A$ is reused across multiple protocol invocations with the same public key (e.g., batch signature verification), the prover can cache the table of evaluations $\{A_C(i,j,k)\}_{i,j,k \in \bits^{\tau+\mu+\nu}}$ and reuse it across invocations, amortizing the NTT cost. Concretely, the prover's online cost for computing Sumcheck~1 round polynomials is dominated by the $O(NMn)$ summation, which does not change; the savings come from avoiding redundant CRT transforms of~$A$.

\subsection{Optimized inner product variant}
\label{sec:inner-product-piop}

When $N=1$, the matrix--vector product $As = t$ reduces to a single inner product $\langle a, s \rangle = t$ in $\ring_p$, where $a \in \ring_p^M$ and $t \in \ring_p$ is a single ring element. This is the relevant setting for SWIFFT hash verification and for proving knowledge of a Module-SIS preimage with a single row. We describe how \cref{prot:ringmul} simplifies in this case.

\paragraph{Structural simplification.}
Setting $N = 1$ gives $\tau = \log N = 0$, so the variable block $\bm{I}$ is empty. The polynomial $\tilde{A}(j, \ell)$ has $\mu + \nu$ variables, $\tilde{s}(j, \ell)$ has $\mu + \nu$ variables, and $\tilde{t}(\ell)$ has $\nu$ variables. The CRT-domain virtual polynomials become
\[
    A_C(j,k) = \sum_{\ell \in \bits^\nu} \tilde{C}(k,\ell)\cdot \tilde{A}(j,\ell), \qquad
    s_C(j,k) = \sum_{\ell \in \bits^\nu} \tilde{C}(k,\ell)\cdot \tilde{s}(j,\ell), \qquad
    t_C(k) = \sum_{\ell \in \bits^\nu} \tilde{C}(k,\ell)\cdot \tilde{t}(\ell).
\]

Sumcheck~1 runs over $\bm{J}$ only ($\mu$ rounds), proving
\[
    e_t = \sum_{j \in \bits^\mu} A_C(j, \alpha') \cdot s_C(j, \alpha')
\]
for verifier challenges $\alpha' \randpick \extensionfield^\nu$, and produces the reduced claim $\sigma' = A_C(\beta, \alpha') \cdot s_C(\beta, \alpha')$ with $\beta \randpick \extensionfield^\mu$.

\paragraph{Holographic $A$, public $t$.}
In the common case where $A$ is public and $t$ is public (or $t = 0$ as in SIS), the protocol reaches its simplest form. The indexer provides $\oracle{A_C}$, and the verifier computes $e_t = t_C(\alpha')$ in $O(n)$ time via the factored evaluation of \cref{lem:m-eval} applied to the known coefficients of~$t$. Concretely, the verifier evaluates
\[
    e_t = \sum_{\ell \in \bits^\nu} \tilde{C}(\alpha', \ell) \cdot \tilde{t}(\ell) = \sum_{k \in \bits^\nu} \meq(\alpha', k) \cdot t_C(k)
\]
by first computing the CRT of~$t$ (an $n$-point NTT, or looked up from the indexed preprocessing) and then taking the $\meq$-weighted inner product with $\alpha'$.

The protocol now involves a single prover oracle $\oracle{\tilde{s}} \in \extensionfield^{\multilin}[\bm{J}, \bm{L}]$ of dimension $\mu + \nu$:
\begin{enumerate}
    \item The verifier samples $\alpha' \randpick \extensionfield^\nu$ and computes $e_t = t_C(\alpha')$.
    \item Sumcheck~1 ($\mu$ rounds): proves $e_t = \sum_{j \in \bits^\mu} A_C(j, \alpha') \cdot s_C(j, \alpha')$, yielding $\sigma' = A_C(\beta, \alpha') \cdot s_C(\beta, \alpha')$.
    \item The verifier queries $e_A \gets \oracle{A_C(\beta, \alpha')}$ and computes $e_s \gets \sigma' / e_A$.
    \item Sumcheck~2 ($\nu$ rounds): proves $e_s = \sum_{\ell \in \bits^\nu} \tilde{C}(\alpha', \ell) \cdot \tilde{s}(\beta, \ell)$, yielding $\sigma'' = \tilde{C}(\alpha', \gamma) \cdot \tilde{s}(\beta, \gamma)$.
    \item The verifier queries $e'_s \gets \oracle{\tilde{s}(\beta, \gamma)}$ and checks $\sigma'' = \tilde{C}(\alpha', \gamma) \cdot e'_s$.
\end{enumerate}
The total proof consists of $\mu + \nu$ sumcheck rounds and a single PCS evaluation of $\tilde{s}$ at the point $(\beta, \gamma) \in \extensionfield^{\mu + \nu}$. This clean structure --- one oracle, one evaluation claim --- composes particularly well with DP24 packing (\cref{sec:dp24-packing}), since there is no need for commitment merging or IOR reductions before invoking the PCS.

\subsection{Reducing the number of commitments}
\label{sec:reducing-commitments}

When the prover commits to multiple oracles, each oracle incurs a Merkle tree and associated PCS opening costs. We describe a standard selector-polynomial technique that merges the oracles for $\tilde{s}$ and $\tilde{t}$ into a single oracle~$h'$, and analyze when this merge is beneficial in conjunction with the IOR and packing optimizations of later sections.

\paragraph{The selector polynomial.}
Define $h' \in \extensionfield^{\multilin}[B, \bm{X}, \bm{L}]$ on $1 + \max(\tau, \mu) + \nu$ variables by
\[
    h'(b, x, \ell) \;:=\; (1 - b)\cdot \tilde{s}(x_{1:\mu}, \ell) \;+\; b\cdot \tilde{t}(x_{1:\tau}, \ell),
\]
where $x_{1:r}$ denotes the first $r$ coordinates of $x$. When $\tau < \mu$, the polynomial $\tilde{t}(x_{1:\tau}, \ell)$ does not depend on coordinates $x_{\tau+1}, \ldots, x_{\mu}$, so $h'$ is well-defined and multilinear. By construction, $h'(0, x, \ell) = \tilde{s}(x_{1:\mu}, \ell)$ and $h'(1, x, \ell) = \tilde{t}(x_{1:\tau}, \ell)$, so a single committed oracle $\oracle{h'}$ replaces both $\oracle{\tilde{s}}$ and $\oracle{\tilde{t}}$.

After Sumcheck~2 (with holographic $A$), the verifier needs the evaluations $\tilde{s}(\beta, \gamma)$ and $\tilde{t}(\alpha, \gamma)$, which correspond to
\begin{align*}
    h'(0,\; \beta,\; \gamma) &= \tilde{s}(\beta, \gamma), \\
    h'(1,\; (\alpha, \beta_{\tau+1:\mu}),\; \gamma) &= \tilde{t}(\alpha, \gamma),
\end{align*}
where we pad $\alpha \in \extensionfield^\tau$ with the coordinates $\beta_{\tau+1}, \ldots, \beta_\mu$ (or any fixed values) to obtain a point of length $\mu$. This gives two evaluation claims on $h'$ at points that differ in the first coordinate ($b = 0$ vs $b = 1$) and possibly in coordinates $1$ through $\tau$ ($\beta_{1:\tau}$ vs $\alpha$), but agree on coordinates $\tau+1$ through $\mu+\nu$ (both use $\beta_{\tau+1:\mu}$ and $\gamma$). These two claims are reduced to a single evaluation via the IOR of \cref{sec:iors}.

\paragraph{When to use $h'$.}
The merge is not always beneficial. Without $h'$, the prover commits to two oracles: $\tilde{s}$ of dimension $\mu + \nu$ and $\tilde{t}$ of dimension $\tau + \nu$. With $h'$, a single oracle of dimension $1 + \mu + \nu$ (assuming $\mu \geq \tau$) replaces both, but the IOR adds a small number of extra field elements to the transcript (see \cref{sec:iors}). The following guidelines cover the main cases:
\begin{itemize}
    \item \emph{$N = 1$, $t$ public (e.g., SIS).} Only $\tilde{s}$ is committed. Do not use $h'$.
    \item \emph{$N = 1$, $t$ private.} Here $\tilde{t}$ has only $\nu$ variables while $\tilde{s}$ has $\mu + \nu$. The merge produces an oracle of dimension $1 + \mu + \nu$, saving one Merkle tree at the cost of one extra PCS round. This is beneficial when $\mu \geq 3$, which holds in all practical MSIS parameter sets. When $\mu \leq 2$, the overhead of the extra variable and IOR outweighs the Merkle tree savings; committing separately is preferable.
    \item \emph{General $N, M$ with $\tau \leq \mu$.} The merge saves the $\tilde{t}$ commitment entirely (one fewer Merkle tree root and PCS opening proof), at the cost of one additional variable in the merged oracle. The net savings in proof size scale with $\min(\tau, \mu) + \nu$: roughly, eliminating the $\tilde{t}$ PCS proof saves $\approx 3(\tau + \nu)$ field elements from PCS folding, while the extra variable adds $\approx 3$ field elements. Thus $h'$ is beneficial whenever $\tau + \nu \geq 2$, which is the generic case.
    \item \emph{$\tau = \mu$ (square matrices).} The two oracles have equal dimension. The merge is always worthwhile, saving one full PCS opening at the cost of one extra round.
\end{itemize}

\subsection{Interactive oracle reductions}
\label{sec:iors}

After the two sumchecks and any holographic oracle queries, the verifier holds one or more evaluation claims on the prover's committed polynomial(s). Before invoking the PCS, we apply interactive oracle reductions (IORs) to consolidate these claims. We follow the framework of Baweja, Chiesa, and Spooner~\cite{BCS24}\psiv{fix cite key}, who give IORs for linear-time-encodable codes in the IOPP model; we describe the two reductions we use and verify their compatibility with random foldable codes (in the BaseFold setting) and with DP24 packing.

\paragraph{Multi-point to single-point IOR.}
Given two evaluation claims $h'(z_0) = v_0$ and $h'(z_1) = v_1$ on the same multilinear polynomial $h'$ of dimension $d$, the standard reduction proceeds as follows. The verifier sends a challenge $r \randpick \extensionfield$ and both parties compute the convex combination $z^* = (1-r)\cdot z_0 + r\cdot z_1$ and $v^* = (1-r)\cdot v_0 + r\cdot v_1$. Since $h'$ is multilinear, the univariate restriction $g(t) := h'((1-t)\cdot z_0 + t\cdot z_1)$ has degree at most $d$ in $t$, and satisfies $g(0) = v_0$, $g(1) = v_1$. The prover sends the coefficients of $g$ (at most $d-1$ additional field elements, since $g(0)$ and $g(1)$ are known), and the verifier checks consistency and accepts the reduced claim $h'(z^*) = g(r)$.

In the special case arising from the $h'$ construction of \cref{sec:reducing-commitments}, the two evaluation points $z_0 = (0, \beta, \gamma)$ and $z_1 = (1, (\alpha, \beta_{\tau+1:\mu}), \gamma)$ differ in at most $1 + \tau$ coordinates (the selector bit and the first $\tau$ entries of $\bm{X}$). When $\tau = 0$ (the $N = 1$ case), the points differ in the first coordinate only, and the restriction $g$ is linear: $g(t) = (1-t)v_0 + t\cdot v_1$, so the prover sends nothing and the verifier simply evaluates $g(r)$. The IOR is free.

For general $\tau$, the prover sends $\tau$ field elements (the intermediate coefficients of $g$), and the reduced claim is $h'(z^*) = g(r)$ at a single point $z^* \in \extensionfield^{1+\mu+\nu}$.

\paragraph{Multi-polynomial to single-polynomial IOR.}
When $h'$ is not used and the prover commits to $\tilde{s}$ and $\tilde{t}$ separately, the verifier holds evaluation claims on distinct polynomials at distinct points: $\tilde{s}(\beta, \gamma) = e'_s$ and $\tilde{t}(\alpha, \gamma) = e'_t$. To reduce these to a single PCS invocation, we first pad both polynomials to a common dimension $d^* = \max(\tau, \mu) + \nu$ by treating the lower-dimensional polynomial as constant in the extra variables. The verifier sends a batching challenge $\zeta \randpick \extensionfield$ and defines the combined polynomial $f = \tilde{s} + \zeta \cdot \tilde{t}_{\mathrm{pad}}$. The two evaluation claims become a single claim on $f$ at a single point (obtained by a further multi-point reduction if the evaluation points differ). The prover commits to $f$ implicitly by providing the combined codeword; the PCS then opens $f$ at the required point.

When additionally $A$ is not holographic and $\oracle{\tilde{A}}$ is prover-committed, its evaluation claim $\tilde{A}(\alpha, \beta, \gamma) = e'_A$ is absorbed into the same batching: $f = \tilde{s} + \zeta \cdot \tilde{t}_{\mathrm{pad}} + \zeta^2 \cdot \tilde{A}$.

\paragraph{Compatibility with random foldable codes and DP24.}
Both IOR reductions operate at the PIOP level, producing a single evaluation claim before the PCS is invoked. The BaseFold PCS then proves this claim using the foldable code's proximity test. The critical property is that the IOR does not impose structural requirements on the code beyond the ability to evaluate the committed polynomial at an arbitrary extension-field point, which random foldable codes support by construction (the encoding is $\field_p$-linear, and evaluation at $\extensionfield$-points is obtained by $\extensionfield$-linear extension of the codeword).

DP24 packing (\cref{sec:dp24-packing}) operates \emph{after} the IOR: the single evaluation claim $f(z^*) = v^*$ on a polynomial of dimension $d^*$ is reduced by the packing sumcheck to a claim on a packed polynomial $\hat{f}$ of dimension $d^* - s$, where $s = \log_2 \kappa$. The IOR's output is a standard evaluation claim, which is exactly the input format that the packing reduction expects. There is no interaction between the IOR mechanics and the packing sumcheck's tower-based decomposition, so compatibility is immediate.

\subsection{DP24 packing}
\label{sec:dp24-packing}

We apply the packing technique of Diamond and Posen~\cite{DP24}\psiv{fix cite key} to reduce the PCS oracle dimension by $s = \log_2 \kappa$ variables, where $\kappa = [\extensionfield : \field_p]$ is the extension degree. This reduces the depth of the Merkle tree and the number of PCS folding rounds, yielding significant concrete proof-size savings. We describe the mechanism, its dependencies, and --- crucially --- why the CRT consistency check (Sumcheck~2) must remain an explicit sumcheck rather than being absorbed into the PCS when packing is in use.

\paragraph{Packing overview.}
Let $f \in \extensionfield^{\multilin}[X_1, \ldots, X_R]$ be the multilinear polynomial to be opened at a point $z \in \extensionfield^R$ (this is the single evaluation claim produced by the IOR of \cref{sec:iors}). The DP24 packing partitions the $R$ variables into $s$ ``packed'' variables $X_1, \ldots, X_s$ and $R - s$ ``free'' variables $X_{s+1}, \ldots, X_R$. Fixing a $\field_p$-basis $\{e_\rho\}_{\rho=1}^\kappa$ of $\extensionfield$ with dual basis $\{e^*_\rho\}$, define the packed polynomial $\hat{f} \in \extensionfield^{\multilin}[X_{s+1}, \ldots, X_R]$ by
\[
    \hat{f}(x_{s+1}, \ldots, x_R) := \sum_{b \in \bits^s} \meq(z_{1:s}, b) \cdot f(b, x_{s+1}, \ldots, x_R).
\]
The evaluation $f(z) = \hat{f}(z_{s+1:R})$ follows by multilinearity. The prover commits to $\hat{f}$ (an oracle of dimension $R - s$) and proves $\hat{f}(z_{s+1:R}) = f(z)$ via the BaseFold PCS using a random foldable code over $\extensionfield$.

The packing reduction itself is an $s$-round sumcheck that verifies the relationship between $f$ and $\hat{f}$. Each round ``lifts'' one tower level, exploiting the tower of quadratic extensions $\field_p \subset \field_{p^2} \subset \cdots \subset \field_{p^\kappa} = \extensionfield$: at level $i$, the sumcheck challenge lives in $\field_{p^{2^i}}$, and the round polynomial has degree~1 in the current variable. After $s$ rounds, the claim reduces to a single evaluation of $\hat{f}$ at an extension-field point, which the BaseFold PCS proves.

\paragraph{Dependencies.}
The packing reduction requires:
\begin{itemize}
    \item A \emph{tower of quadratic extensions} $\field_p \subset \field_{p^2} \subset \cdots \subset \extensionfield$ with explicit embedding and projection maps. For odd primes~$p$, we construct such a tower in \cref{sec:odd-prime-tower}\psiv{fix cref}.
    \item \emph{Linear-time encodable random foldable codes} over each tower level $\field_{p^{2^i}}$, $i = 0, \ldots, s$. We describe our construction based on~\cite{BaseFold}\psiv{fix cite} in \cref{sec:foldable-codes}\psiv{fix cref}.
\end{itemize}

\paragraph{Transcript costs.}
The $s$-round packing sumcheck contributes $3s$ field elements to the PIOP transcript (one degree-1 univariate per round, costing 2 coefficients, plus 1 claimed evaluation --- but since the constant term is determined by the claim, each round sends 1 new field element; the total is $2s$\psiv{double-check this count against DP24}). The BaseFold PCS then runs $R - s$ folding rounds instead of $R$, saving $3s$ field elements from PCS round messages. The net transcript-size change is approximately zero, but the Merkle tree depth decreases by $s$, which saves $s \cdot \lambda$ bits of hash proof per query (where $\lambda$ is the hash output length). For $s = 4$ and $\lambda = 256$, this is a saving of 128 bytes per query repetition.

\paragraph{Incompatibility of CRT absorption with packing.}
\label{par:crt-absorption-incompatibility}

A natural question is whether Sumcheck~2 (the CRT consistency check) can be absorbed into the PCS as a weighted-sum constraint, eliminating it as an explicit sumcheck and saving $\sim 3\nu$ field elements. We now show that this absorption is \emph{incompatible} with DP24 packing, and therefore Sumcheck~2 must remain explicit whenever packing is used.

\begin{proposition}\label{prop:crt-packing-incompatibility}
    Let $\hat{f}$ be the packed polynomial obtained from $f$ via the DP24 reduction with packing parameter $s \geq 1$. Let $\varphi_\rho : \extensionfield \to \extensionfield$ denote the $\field_p$-linear dual-basis extraction map $\varphi_\rho(x) = \langle e^*_\rho, x \rangle_{\field_p}$ (i.e., $\varphi_\rho$ extracts the $\rho$-th $\field_p$-coordinate of $x$). If the CRT consistency check is absorbed into the PCS, the resulting constraint on $\hat{f}$ takes the form
    \[
        \sum_{y} w(y, \rho) \cdot \varphi_\rho\!\bigl(\hat{f}(y)\bigr) = \sigma_\rho
    \]
    for verifier-computable weights $w(\cdot)$ and each basis index $\rho$. This constraint is $\field_p$-linear but not $\extensionfield$-linear in the evaluations of $\hat{f}$.

    The BaseFold folding step replaces $\hat{f}$ with $\hat{f}_{\mathrm{new}}(y') = \hat{f}(0, y') + r \cdot \hat{f}(1, y')$ for a challenge $r \in \extensionfield$. For the constraint to be maintained through folding, we would need
    \[
        \varphi_\rho\bigl(\hat{f}(0, y') + r \cdot \hat{f}(1, y')\bigr) = \varphi_\rho\bigl(\hat{f}(0, y')\bigr) + r \cdot \varphi_\rho\bigl(\hat{f}(1, y')\bigr)
    \]
    for all $r \in \extensionfield$ and all evaluations. But $\varphi_\rho$ is $\field_p$-linear, not $\extensionfield$-linear: for general $r \in \extensionfield \setminus \field_p$, $\varphi_\rho(r \cdot x) \neq r \cdot \varphi_\rho(x)$. Hence the absorbed constraint cannot be tracked through $\extensionfield$-challenge folding.
\end{proposition}

\begin{proof}
    It suffices to exhibit a concrete failure. Let $p$ be an odd prime and $\extensionfield = \field_{p^2}$ with basis $\{1, \theta\}$ where $\theta^2 = g$ for a non-square $g \in \field_p$. Then $\varphi_1(a + b\theta) = a$ and $\varphi_2(a + b\theta) = b$. Take $r = \theta$ and $x = 1$. Then $\varphi_1(\theta \cdot 1) = \varphi_1(\theta) = 0$, but $\theta \cdot \varphi_1(1) = \theta \cdot 1 = \theta \neq 0$. The constraint breaks at the first folding step.
\end{proof}

By contrast, standard evaluation claims $\hat{f}(z) = v$ are trivially compatible with packing: the verifier holds $v$ and applies $\varphi_\rho$ post-hoc if needed, without requiring the PCS to track an $\field_p$-linear constraint through its folds.

\paragraph{Correct protocol architecture.}
The incompatibility above dictates the following structure when all optimizations are combined:
\begin{enumerate}
    \item \textbf{Sumcheck~1} ($\mu$ rounds, or $\tau + \mu$ rounds if $\bm{I}$ is nonempty): the matrix-vector product check in the CRT domain, producing evaluation claims on $A_C$, $\tilde{s}$, and $\tilde{t}$.
    \item \textbf{Sumcheck~2} ($\nu$ rounds): the explicit CRT consistency check, reducing CRT-domain claims to coefficient-domain evaluation claims on the committed oracles. When $A$ is holographic, only $\tilde{s}$ and $\tilde{t}$ (or $\tilde{s}$ alone, if $t$ is also public) participate.
    \item \textbf{IOR reduction} (\cref{sec:iors}): consolidates the resulting evaluation claims into a single claim $f(z) = v$ on a single polynomial of dimension $d^*$. When the $h'$ merge is used, this is a multi-point IOR; otherwise, a multi-polynomial IOR.
    \item \textbf{DP24 packing sumcheck} ($s$ rounds): reduces the claim to $\hat{f}(z') = v$ on a packed polynomial of dimension $d^* - s$.
    \item \textbf{BaseFold PCS} ($d^* - s$ folding rounds): proves the evaluation of $\hat{f}$ using a random foldable code over $\extensionfield$, with linear-time encoding (\cref{sec:foldable-codes}\psiv{fix cref}).
\end{enumerate}
The total PIOP transcript size is $3(\mu + \nu) + O(\tau) + 2s + 3(d^* - s)$ field elements, plus the IOR messages. The prover complexity is $O(NMn)$ for the sumchecks plus $O(2^{d^*})$ for the PCS encoding; the verifier complexity is $O(n + \log NM)$ field operations.

\paragraph{Batch verification across MSIS instances.}
\label{par:batch-verification}
A common setting in lattice-based cryptography is the need to verify $k$ independent Module-SIS relations $A s_i = t_i$ for $i = 1, \ldots, k$ with a shared public matrix~$A$. Since the indexed oracle $\oracle{A_C}$ is the same across instances, the indexing cost is paid once. The verifier batches the $k$ instances by sampling $\eta_1, \ldots, \eta_k \randpick \extensionfield$ and checking the combined relation
\[
    A \Bigl(\sum_{i=1}^k \eta_i s_i\Bigr) = \sum_{i=1}^k \eta_i t_i.
\]
This reduces $k$ protocol executions to a single execution on the combined witness $s^* = \sum_i \eta_i s_i$ and target $t^* = \sum_i \eta_i t_i$, at the cost of the prover computing the random linear combination of witness vectors. If each $s_i$ has dimension $M$, the combination costs $O(kMn)$ ring-element operations. The PCS can additionally batch the $k$ oracle openings using a single multi-polynomial IOR (treating each $\tilde{s}_i$ as a separate polynomial), or the prover can commit to $\tilde{s}^*$ directly if the combination is performed before commitment. The latter is preferable when $k$ is known in advance (e.g., batch signature verification), as it requires a single oracle of the same dimension as each individual $\tilde{s}_i$.